\documentclass[11pt]{article}
\usepackage[margin=1in]{geometry}
\usepackage{amsmath, amssymb}
\usepackage{booktabs}
\usepackage{hyperref}

\title{Theoretical Minimum $\sigma_\mathrm{TOF}$ from Correlation Peak FWHM and SNR}
\author{BATPING Project Technical Notes}
\date{July 2026}

\begin{document}
\maketitle

\section{Setup}

The time-of-flight (TOF) is estimated from the peak of the Hilbert envelope of the
autocorrelation of the received signal.  Near the peak, the envelope is well approximated
by a Gaussian:
\begin{equation}
  E(\tau) \;=\; A_\mathrm{TOF}\,\exp\!\left(-\frac{(\tau-\tau_0)^2}{2\sigma_\mathrm{env}^2}\right)
             \;+\; n(\tau),
  \label{eq:envelope}
\end{equation}
where $A_\mathrm{TOF}$ is the peak amplitude, $\tau_0$ is the true TOF,
$n(\tau)$ is additive white noise with RMS $\sigma_n$, and $\sigma_\mathrm{env}$
is the width of the envelope.  The FWHM of the peak relates to $\sigma_\mathrm{env}$ by
\begin{equation}
  \mathrm{FWHM} \;=\; 2\sqrt{2\ln 2}\;\sigma_\mathrm{env} \;\approx\; 2.355\,\sigma_\mathrm{env}.
  \label{eq:fwhm}
\end{equation}

\section{Cramér--Rao Bound}

The Fisher information for the location parameter $\tau_0$ of a Gaussian peak
in white noise is
\begin{equation}
  \mathcal{I}(\tau_0)
    \;=\; \frac{1}{\sigma_n^2}\int_{-\infty}^{\infty}
          \left[\frac{\partial E}{\partial \tau_0}\right]^2 d\tau.
\end{equation}
Computing the derivative from \eqref{eq:envelope}:
\begin{equation}
  \frac{\partial E}{\partial \tau_0}
    \;=\; A_\mathrm{TOF}\,\frac{\tau - \tau_0}{\sigma_\mathrm{env}^2}
          \exp\!\left(-\frac{(\tau-\tau_0)^2}{2\sigma_\mathrm{env}^2}\right).
\end{equation}
Substituting and evaluating the Gaussian integral:
\begin{equation}
  \mathcal{I}(\tau_0)
    \;=\; \frac{A_\mathrm{TOF}^2}{\sigma_n^2\,\sigma_\mathrm{env}^4}
          \int_{-\infty}^{\infty}
          (\tau-\tau_0)^2\,
          e^{-(\tau-\tau_0)^2/\sigma_\mathrm{env}^2}\,d\tau
    \;=\; \frac{A_\mathrm{TOF}^2}{\sigma_n^2\,\sigma_\mathrm{env}^2}
          \cdot\frac{\sqrt{\pi}}{2}\cdot\frac{1}{\sqrt{1/\sigma_\mathrm{env}^2}}.
\end{equation}
After simplification:
\begin{equation}
  \mathcal{I}(\tau_0)
    \;=\; \frac{A_\mathrm{TOF}^2\,\sqrt{\pi\,\sigma_\mathrm{env}^2}}
               {2\,\sigma_n^2\,\sigma_\mathrm{env}^2}
    \;=\; \frac{A_\mathrm{TOF}^2\sqrt{\pi}}{2\,\sigma_n^2\,\sigma_\mathrm{env}}.
\end{equation}
The Cramér--Rao bound gives the minimum variance of any unbiased estimator:
\begin{equation}
  \sigma_\tau^2 \;\geq\; \frac{1}{\mathcal{I}(\tau_0)}
    \;=\; \frac{2\,\sigma_n^2\,\sigma_\mathrm{env}}{A_\mathrm{TOF}^2\,\sqrt{\pi}}.
\end{equation}
Defining the \emph{amplitude SNR} at the TOF peak as
$\mathrm{SNR}_A \equiv A_\mathrm{TOF}/\sigma_n$, and using
$\sigma_\mathrm{env} = \mathrm{FWHM}/2.355$:
\begin{equation}
  \sigma_\tau^2 \;\geq\;
    \frac{2\,\sigma_\mathrm{env}}{\mathrm{SNR}_A^2\,\sqrt{\pi}}
    \;=\;
    \frac{2}{\sqrt{\pi}}\cdot\frac{\mathrm{FWHM}}{2.355\,\mathrm{SNR}_A^2}.
\end{equation}
Taking the square root and absorbing numerical constants:
\begin{equation}
  \boxed{
    \sigma_\tau^{\mathrm{theory}}
      \;\approx\;
      \frac{\mathrm{FWHM}}{2.355\,\sqrt{2}\;\mathrm{SNR}_A}
      \;=\;
      \frac{\mathrm{FWHM}}{3.33\;\mathrm{SNR}_A}.
  }
  \label{eq:main}
\end{equation}

\section{Definitions of Measured Quantities}

\paragraph{FWHM.}
Measured from the Hilbert envelope of the autocorrelation at the TOF peak,
using parabolic sub-sample interpolation at the half-maximum crossings.

\paragraph{$\mathrm{SNR}_A$.}
The amplitude SNR at the TOF peak:
\begin{equation}
  \mathrm{SNR}_A \;=\; \frac{A_\mathrm{TOF}}{\sigma_n},
\end{equation}
where $A_\mathrm{TOF}$ is the Hilbert envelope value at the detected TOF lag, and
$\sigma_n$ is the RMS of the Hilbert envelope in a noise-only region well away from
any acoustic arrival (in practice, lags $15$--$20\,\mathrm{ms}$, where no signal energy
is expected).

In the normalized autocorrelation (lag-0 peak $= 1$), both $A_\mathrm{TOF}$ and
$\sigma_n$ are expressed as fractions of the lag-0 peak.

\paragraph{Temperature sensitivity.}
The speed of sound in air is $c(T) = 331.3 + 0.606\,T$ m/s.  For a round-trip path
of length $2L$, the TOF sensitivity to temperature is
\begin{equation}
  \frac{\partial \tau}{\partial T}
    \;=\; -\frac{2L}{c^2}\,\frac{\partial c}{\partial T}
    \;\approx\; -10.3\;\mu\mathrm{s/{}^\circ C}
\end{equation}
at $T = 20\,^\circ$C with $L = 1\,\mathrm{m}$.  The theoretical temperature precision is
therefore
\begin{equation}
  \sigma_T^{\mathrm{theory}}
    \;=\; \frac{\sigma_\tau^{\mathrm{theory}}}{10.3\;\mu\mathrm{s/{}^\circ C}}.
\end{equation}

\section{Numerical Example}

For the best result from the July 2 recordings
(13\,kHz compensated chirp, 10\,ms duration):
\begin{align*}
  \mathrm{FWHM} &= 99\;\mu\mathrm{s}, \\
  \sigma_\tau^{\mathrm{measured}} &= 0.34\;\mu\mathrm{s}, \\
  \mathrm{SNR}_A^{\mathrm{implied}} &= \frac{99}{3.33 \times 0.34} \approx 87
    \quad(\approx 39\,\mathrm{dB}).
\end{align*}
Once $\sigma_n$ is measured directly from the noise floor of the autocorrelation,
equation~\eqref{eq:main} can be evaluated forward (predicting $\sigma_\tau$) rather
than backward (inferring $\mathrm{SNR}_A$ from the measured $\sigma_\tau$).

\section{Interpretation}

Equation~\eqref{eq:main} makes explicit the two independent routes to improving
$\sigma_\tau$:

\begin{enumerate}
  \item \textbf{Narrow the FWHM} by increasing the effective signal bandwidth ---
        through wider chirp passbands, amplitude compensation, or higher-order
        modulation schemes.
  \item \textbf{Raise the SNR} by increasing drive voltage, using co-added
        templates (SNR $\propto \sqrt{N}$), reducing ambient noise, or improving
        transducer coupling efficiency.
\end{enumerate}
The two are not fully independent: a wider passband driven by compensation often
pushes into lower-SNR frequency regions, so gains in FWHM may come at a cost to
$\mathrm{SNR}_A$.  The product $\mathrm{SNR}_A \times \mathrm{FWHM}^{-1}$ is the
figure of merit that directly determines $\sigma_\tau$.

\end{document}
